\documentclass[10pt,a4paper]{amsart}
\usepackage[margin=19mm]{geometry}
\usepackage{amsmath,amssymb,mathtools}
\usepackage[scr=rsfs,cal=euler]{mathalfa}
\usepackage{xfrac}
\usepackage[unicode,hidelinks,bookmarks=false]{hyperref}
\newcommand{\ie}{i.e.\ }
\newcommand{\cf}{cf.\ }
\newcommand{\resp}{respectively\ }
\newcommand{\Subsection}{\subsection}
\providecommand{\nobreakdash}{\nobreak\hbox{-}}
\setlength{\emergencystretch}{2em}
\multlinegap0pt
\usepackage{bm}
\usepackage{bbm}
\usepackage{dsfont}
\usepackage{xspace}
\usepackage{cancel}
\usepackage{mathtools}
\usepackage[shortlabels]{enumitem}
\usepackage{amsthm}
\makeatletter
\@ifundefined{theorem}{\newtheorem{theorem}{Theorem}}{}
\@ifundefined{proposition}{\newtheorem{proposition}[theorem]{Proposition}}{}
\makeatother
\newcommand{\K}{{\mathcal{K}}}  
\newcommand{\bC}{{\mathbb{C}}}
\newcommand{\bF}{{\mathbb{F}}}
\newcommand{\bM}{{\mathbb{M}}}
\newcommand{\bT}{{\mathbb{T}}}
\newcommand{\bZ}{{\mathbb{Z}}}
\newcommand{\rC}{\mathrm{C}}
\renewcommand{\H}{{\mathcal{H}}}
\renewcommand{\phi}{\varphi}
\title[On the universal commuting dilation constant]{On the universal commuting dilation constant}
\author{Ian Thompson}
\date{Source: arXiv:2606.12506v1 (CC BY 4.0)}
\begin{document}
\maketitle

\noindent\emph{Source and licence.} On the universal commuting dilation constant. Version arXiv:2606.12506v1 (CC BY 4.0), distributed under the Creative Commons Attribution 4.0 International licence, as recorded in the arXiv licence metadata for this exact version and shown on the abstract page https://arxiv.org/abs/2606.12506v1. Full rights notice: licenses/arxiv-2606.12506-CC-BY-4.0.txt.\par\smallskip

\noindent\textit{Editorial scope.} Counted unit: the Theorem whose source label is \texttt{t:uni-comm-dil-constant} in the article (source0.tex lines 465--467, source label \texttt{t:uni-comm-dil-constant}), delivered together with the article's own complete original proof of it (source0.tex lines 469--526, one \texttt{proof} environment of 58 lines). No mathematics is written, repaired or re-derived: statement and proof are the article's own text line for line. Required context retained uncounted: p:cocyc-uni-comm-uni (lines 390--396); p:univ-uni-compress (lines 414--427). Every definition, display and label of those lines is retained unchanged. Omitted from this package and not redistributed: 1--389, 397--413, 428--464, 468--468, 527--1075. Nothing inside a retained excerpt is omitted or altered: every retained body is delivered exactly as the article's lines are written, so no omission receipt of any kind accompanies this package. Citations: the retained excerpts carry no citation command at all, so no bibliography entry of the article is transcribed and no reference is redistributed. Reference adaptation: none was needed and none was made; every reference inside the delivered unit resolves inside it, so no unresolved reference or citation is produced. Printed slips kept literal: no printed word, symbol, hypothesis or number of the counted statement or of its proof was altered. Layout and class: the article's own class is \texttt{amsart} with packages including graphicx, amscd, amsmath, amsthm, amsfonts, amssymb, mathrsfs, bm, enumerate, amsrefs, xcolor, hyperref, comment, tikz; the wrapper is the lane's pinned \texttt{amsart} editorial wrapper at 10pt on A4, which restates the theorem-like environments the retained text needs and adds the article's own macro definitions that the retained text needs (\texttt{\textbackslash H}, \texttt{\textbackslash K}, \texttt{\textbackslash bC}, \texttt{\textbackslash bF}, \texttt{\textbackslash bM}, \texttt{\textbackslash bT}, \texttt{\textbackslash bZ}, \texttt{\textbackslash phi}, \texttt{\textbackslash rC}), with extra wrapper packages (bm, bbm, dsfont, xspace, cancel, mathtools, enumitem). The delivered numbering is the unit's own. Assets: no retained span contains a figure environment, a graphics-inclusion command, a PDF-page inclusion, a table or a TikZ picture; no image file is copied into this package, the package declares no asset dependency and no image file is redistributed with it. Licence: version arXiv:2606.12506v1 of this article is distributed under the Creative Commons Attribution 4.0 International licence, as recorded in the arXiv licence metadata for this exact version and shown on the abstract page https://arxiv.org/abs/2606.12506v1. A full rights notice is delivered at licenses/arxiv-2606.12506-CC-BY-4.0.txt. Characters: every retained excerpt is pure ASCII, so the delivered engine, XeLaTeX with its default Latin Modern text font, reports no missing character.
\begin{proposition}\label{p:cocyc-uni-comm-uni}
Let $\K$ be a Hilbert space and $S\in B(\K)$ be a unitary operator. Assume that $\rC^*(\bF_2)\subseteq B(\H)$ for some Hilbert space $\H$ and consider the pair of operators $Z,W\in B(\H\otimes \ell^2(\bZ)\otimes\K)$ defined by
    \[
        Z(h\otimes\delta_m\otimes\eta) = \mathfrak{u}_ah\otimes\delta_{m+1}\otimes\eta, \qquad W(h\otimes\delta_m\otimes\eta) = \mathfrak{u}_{a^m ba^{-m}}h\otimes\delta_m\otimes S\eta,
    \]
for $h\in\H$, $m\in\bZ$, and $\eta\in\K$. Then, $Z$ and $W$ are commuting unitary operators.
\end{proposition}

\begin{proposition}\label{p:univ-uni-compress}
    Let $\K$ be a Hilbert space and $S\in B(\K)$ be some unitary operator. Assume that $\rC^*(\bF_2)\subseteq B(\H)$ for some Hilbert space $\H$ and consider the pair of commuting unitaries $Z,W\in B(\H\otimes \ell^2(\bZ)\otimes\K)$ defined by
    \[
        Z(h\otimes\delta_m\otimes\eta) = \mathfrak{u}_ah\otimes\delta_{m+1}\otimes\eta, \qquad W(h\otimes\delta_m\otimes\eta) = \mathfrak{u}_{a^m ba^{-m}}h\otimes\delta_m\otimes S\eta,
    \]
    for $h\in\H, m\in\bZ$, and $\eta\in K$. Assume that $\xi\in\ell^2(\bZ)\otimes\K$ is a unit vector given by $\xi=\sum_{m\in\bZ}\delta_m\otimes\xi_m$ for some sequence $(\xi_m)_m$ in $\K$, and define an isometry $V:\H\rightarrow \H\otimes \ell^2(\bZ)\otimes\K$ by $Vh=h\otimes\xi$. Then, the following statements hold.
    \begin{enumerate}[{\rm (i)}]
        \item We have that
        \[
            V^*ZV = \left(\sum_{m\in\bZ}\langle \xi_m,\xi_{m+1}\rangle \right)\mathfrak{u}_a \quad \text{and} \quad V^*WV = \sum_{m\in\bZ} \langle S\xi_m,\xi_m\rangle\mathfrak{u}_{a^m b a^{-m}}.
        \]
        \item We have that $V^* \rC^*(Z,W) V \subseteq \rC^*(\bF_2)$.
    \end{enumerate}
\end{proposition}

\begin{theorem}\label{t:uni-comm-dil-constant}
    There is a unital completely positive map $\Phi:\rC(\bT^2)\rightarrow\rC^*(\bF_2)$ such that $\Phi(z) = \frac{\sqrt{\phi}}{2}\mathfrak{u}_a$ and $\Phi(w) = \frac{\sqrt{\phi}}{2}\mathfrak{u}_b$.
\end{theorem}

\begin{proof}
    Throughout, assume that $\rC^*(\bF_2)\subseteq B(\H)$ for some Hilbert space $\H$. Let $e_1, e_2$ be the standard orthonormal basis of $\bC^2$ and consider the unitary operator $S\in \bM_2$ defined by $Se_1=e_1$ and $Se_2=-e_2$.
    
    First, we construct an appropriate unit vector $\xi\in\ell^2(\bZ)\otimes\bC^2$ of the form $\sum_{m\in\bZ} \delta_m\otimes\xi_m$ for which to apply Proposition \ref{p:univ-uni-compress}. To this end, we consider the following constants
    \[
        p=\frac{5+\sqrt5}{10} \quad \text{and} \quad r=\sqrt{\frac{2(\sqrt5-1)}{5}},
    \]
    and note that $0<r<p<1$. Now, set
    \[
        \xi_0 = \sqrt{\frac{p+r}{2}}e_1 + \sqrt{\frac{p-r}{2}}e_2,
    \]
    and observe that $\|\xi_0\|^2 = p$. Furthermore, define
    \[
        q=\frac{1-p}{2}, \quad \rho=\sqrt{\frac{q}{\frac{1}{2}(p+\sqrt{p^2-r^2})+q}}, \quad \text{and} \quad \alpha_m = \sqrt{q(1-\rho^2)}\,\rho^{m-1}, \ \ \ m\geq 1.
    \]
    Then, we consider
    \[
        \xi=\sum_{m\in\bZ}\delta_m\otimes \xi_m\in \ell^2(\bZ)\otimes\bC^2
    \]
    where $\xi_m = \alpha_{|m|}\frac{e_1+e_2}{\sqrt2}$ for each $m\neq 0$. Note that $\xi$ is a unit vector since
    \[
        \|\xi\|^2 = \|\xi_0\|^2 +2\sum_{m\geq1}\alpha_m^2 = p + 2q(1-\rho^2)\sum_{m\geq1}\rho^{2(m-1)} = p+2q = 1.
    \]

    By Proposition \ref{p:cocyc-uni-comm-uni}, the pair of operators on $\H\otimes\ell^2(\bZ)\otimes\bC^2$ defined by
    \[
        Z(h\otimes\delta_m\otimes\eta) = \mathfrak{u}_ah\otimes\delta_{m+1}\otimes\eta \quad \text{and} \quad W(h\otimes\delta_m\otimes\eta) = \mathfrak{u}_{a^mba^{-m}}h\otimes\delta_m\otimes S\eta
    \]
    are commuting unitaries. Thus, there is a unital $*$-representation $\pi:\rC(\bT^2)\rightarrow B(\H\otimes\ell^2(\bZ)\otimes\bC^2)$ such that $\pi(z) = Z$ and $\pi(w) = W$. Let $V:\H\rightarrow\H\otimes\ell^2(\bZ)\otimes\bC^2$ be the isometry given by $Vh = h\otimes\xi$ and define a unital completely positive map
    \[
        \Phi_1:\rC(\bT^2)\longrightarrow B(\H), \qquad f\longmapsto V^*\pi(f)V.
    \]
    By Proposition \ref{p:univ-uni-compress}, we have that $\Phi_1(\rC(\bT^2))\subseteq\rC^*(\bF_2)$ and that
    \[
        \Phi_1(z) = \left( \sum_{m\in\mathbb Z} \langle \xi_m,\xi_{m+1}\rangle\right)\mathfrak{u}_a \quad \text{and} \quad \Phi_1(w) = \sum_{m\in\mathbb Z}\langle S\xi_m,\xi_m\rangle\mathfrak{u}_{a^m b a^{-m}}.
    \]
    For each $m\neq 0$, note that $\langle S\xi_m,\xi_m\rangle = 0$ and so,
    \[
        \Phi_1(w)= \langle S\xi_0,\xi_0\rangle\mathfrak{u}_b =r\mathfrak{u}_b.
    \]
    Setting $A = \langle \xi_0,\frac{e_1+e_2}{\sqrt{2}}\rangle$, we then find that
    \begin{align*}
        \sum_{m\in\bZ}\langle\xi_m, \xi_{m+1}\rangle & =2A\alpha_1  +  2\sum_{m\geq 1}\langle \xi_m, \xi_{m+1}\rangle\\
        & = 2A\alpha_1 + 2 \sum_{m\geq 1}\alpha_m \alpha_{m+1}\\
        & = 2A\sqrt{\frac{A^2q}{A^2+q}} +2q \sqrt{\frac{q}{A^2+q}}\\
        & = 2\sqrt{q\left(A^2 + q \right)} = \sqrt{\frac{1+\sqrt{5}}{10}}.
    \end{align*}
    Therefore, $\Phi_1(z)=\widetilde{r} \mathfrak{u}_a$ where $\widetilde{r}:=\sqrt{\frac{1+\sqrt{5}}{10}}$.
    
    Let $\beta$ be the $*$-automorphism of $\rC^*(\bF_2)$ induced by the group automorphism of $\bF_2$ that exchanges $a$ and $b$. Similarly, let $\sigma$ be the $*$-automorphism of $\rC(\bT^2)$ that exchanges the coordinate functions $z$ and $w$. Define a unital completely positive $\Phi_2:\rC(\bT^2)\rightarrow\rC^*(\bF_2)$ by $\Phi_2 =\beta\circ \Phi_1\circ \sigma$. Furthermore, define a unital completely positive map $\Phi:\rC(\bT^2)\rightarrow\rC^*(\bF_2)$ by $\Phi=\frac{\Phi_1+\Phi_2}{2}$. Then, $\Phi$ satisfies
    \[
        \Phi(z)=\frac{r+\widetilde{r}}{2}\mathfrak{u}_a \quad \text{and} \quad \Phi(w)=\frac{r+\widetilde{r}}{2}\mathfrak{u}_b.
    \]
    Finally, observe that
    \[
        \frac{r+\widetilde{r}}{2} = \frac{1}{2}\left(\sqrt{\frac{2(\sqrt{5}-1)}{5}} + \sqrt{\frac{1+\sqrt{5}}{10}} \right) = \frac{1}{4}\sqrt{2(1+\sqrt{5})} = \frac{\sqrt{\phi}}{2}.
    \]
\end{proof}

\end{document}
